% Informe 2 - Implementación (máximo 15 páginas sin contar referencias).
% Debe incluir las correcciones del informe anterior solicitadas por la
% ayudante y la respuesta a cada aspecto de la Retroalimentación del Estado
% de Avance (indicando qué fue solicitado). Redacción en tercera persona.
\documentclass[letter, 10pt]{article}
\usepackage[utf8]{inputenc}
\usepackage[T1]{fontenc}
\usepackage[spanish]{babel}
\usepackage{amsfonts}
\usepackage{amsmath}
\usepackage{graphicx}
\usepackage{url}
\usepackage[top=3cm,bottom=3cm,left=3.5cm,right=3.5cm,footskip=1.5cm,headheight=1.5cm,headsep=.5cm]{geometry}

\begin{document}
\title{Inteligencia Artificial \\ \begin{Large}Implementación: Problema ``Segmentación de imágenes médicas mediante metaheurísticas''\end{Large}}
\author{Matías Peñaloza 2023730378}
\date{\today}
\maketitle

\begin{abstract}
% TODO
\end{abstract}

\section{Introducción}
% TODO: contexto y motivación~\cite{Otsu1979}.

\section{Definición del problema}
% TODO

\section{Estado del arte}
% TODO: incluir las correcciones del Informe 1.

\section{Modelo matemático}
% TODO: variables, restricciones y función objetivo.

\section{Representación}
% TODO

\section{Descripción del algoritmo}
% TODO: función de evaluación con tablas precalculadas~\cite{Liao2001Fast},
% solución inicial (Greedy), Tabú Search, movimientos, manejo de
% restricciones.

\section{Experimentos}
% TODO: instancias, parámetros, entorno de ejecución.

\section{Resultados}
% TODO: tablas y gráficos comparativos; comparación con investigaciones
% previas~\cite{Hammouche2010Comparative}.

\section{Conclusiones}
% TODO

\section{Uso LLMs}
% TODO: herramienta, propósito, prompts relevantes, nivel de edición y
% validación, de acuerdo con los lineamientos del proyecto.

% TODO: eliminar este \nocite cuando existan citas reales en el texto
% (no deben quedar referencias sin usar).
\nocite{Otsu1979, Liao2001Fast}
\bibliographystyle{plain}
\bibliography{Referencias}

\end{document}
